% ======================================================================
% SECTION 01: EXECUTIVE SUMMARY
% ======================================================================
\section{Executive Summary}
\noindent\textbf{GitHub Repository:} \url{https://github.com/AngeloIbanez/earthquake-monitoring-team_1}\\[0.15cm]
In Workshop~\#1, Team~1 established the geophysical foundations of global earthquake activity monitoring, validated four open seismological feeds (USGS FDSN Historical, EMSC FDSN Historical, USGS Real-Time GeoJSON, and GDACS RSS~2.0 XML), and designed a seven-table canonical relational data model. In Workshop~\#2, we deployed a containerized Apache Airflow~2.10.4 and PostgreSQL~16 environment via Docker Compose, implementing the initial \texttt{bronze\_ingestion\_dag} and FileSensor-triggered \texttt{silver\_processing\_dag}.

In Workshop~\#3, we transition from single-source ingestion to full multi-agency integration and implement the **Gold layer**---the analytical and decision-support tier of the Medallion Architecture~\cite{armbrust2021lakehouse}. Our work fulfills three interconnected engineering objectives:
\begin{enumerate}
    \item \textbf{Bronze \& Silver Pipeline Refinements and Canonical Integration:} We expanded \texttt{bronze\_ingestion\_dag.py} and \texttt{silver\_processing\_dag.py} (alongside the standalone CLI \texttt{ingest.py}) to extract, validate, and persist all four configured seismic sources across the seven canonical PostgreSQL tables (\texttt{source}, \texttt{ingestion\_run}, \texttt{ingestion\_state}, \texttt{event}, \texttt{earthquake}, \texttt{event\_group}, and \texttt{rejected\_record}). We eliminated silent null-to-zero coercions, activated JSONB quarantine logging, implemented incremental watermark overlap windows ($24\,\text{h}$ for USGS, $48\,\text{h}$ for GDACS), and integrated transitive cross-catalog event matching.
    \item \textbf{Distributed Silver-to-Gold Transformation with PySpark:} We engineered a custom Docker image extending \texttt{apache/airflow:2.10.4} with OpenJDK~17 (\texttt{default-jre-headless}) and PySpark~3.5+. Within \texttt{airflow/dags/gold\_processing\_dag.py} and \texttt{src/gold\_pyspark.py}, PySpark runs in \texttt{local[*]} mode to read the entire \texttt{datalake\_silver/} directory as a unified distributed DataFrame (\texttt{spark.read.parquet("datalake\_silver/")}), simulating a cloud data lake folder read pattern.
    \item \textbf{Data Governance KPIs \& Functional User Storytelling Summaries:} Using both the PySpark DataFrame API and Spark SQL, the Gold pipeline computes and persists two timestamped Parquet datasets in \texttt{datalake\_gold/}:
    \begin{itemize}
        \item \texttt{governance\_YYYYMMDD\_HHMMSS.parquet} ($113$ rows $\times$ $13$ columns): Grounded in empirical Silver data quality observations, tracking column null rates, implicit $10.0\,\text{km}$ default depth pinning ($29.85\%$ overall; $98.25\%$ in EMSC), intra-lake duplicate rates ($23.10\%$), cross-source match pairs ($637$ pairs), schema compliance ($100.0\%$), IQR outlier rates, and \texttt{place} string length distributions.
        \item \texttt{storytelling\_YYYYMMDD\_HHMMSS.parquet} ($84$ rows $\times$ $14$ columns): Aligned with Workshop~\#1 functional user stories for civil protection and seismic risk analysts, delivering Gutenberg-Richter magnitude tier distributions, daily radiated seismic energy ($\log_{10} E$), the **Seismic Activity Anomaly Index (SAAI)**, top active tectonic regions, multi-source catalog comparisons, and focal depth regime breakdowns.
    \end{itemize}
\end{enumerate}

\begin{figure*}[!t]
\centering
\resizebox{0.98\textwidth}{!}{%
\begin{tikzpicture}[
    node distance=0.75cm and 1.0cm,
    box/.style={rectangle, draw=black!80, thick, rounded corners=3pt, minimum width=3.0cm, minimum height=1.1cm, align=center, font=\footnotesize\sffamily, fill=white},
    layer/.style={rectangle, draw=navyblue!70, dashed, thick, rounded corners=5pt, inner sep=7pt},
    arrow/.style={-{Stealth[length=2.4mm, width=1.7mm]}, thick, draw=black!80}
]

% 1. Sources
\node[box, fill=blue!8] (s1) {\textbf{USGS FDSN \& RT}\\GeoJSON REST\\(\texttt{usgs\_fdsn / realtime})};
\node[box, fill=blue!8, below=0.35cm of s1] (s2) {\textbf{EMSC \& GDACS}\\JSON REST \& RSS XML\\(\texttt{emsc\_fdsn / gdacs})};

% 2. Bronze
\node[box, fill=orange!15, right=1.15cm of s1] (bdag) {\textbf{Bronze Ingestion DAG}\\(\texttt{bronze\_ingestion\_dag})\\Multi-Source Extract\\Schedule: \texttt{*/15 * * * *}};
\node[box, fill=orange!8, below=0.35cm of bdag] (bstore) {\textbf{Bronze Landing Zone}\\\texttt{datalake\_bronze/}\\\texttt{earthquakes\_*.json}\\\texttt{emsc\_*.json, gdacs\_*.xml}};

% 3. Silver
\node[box, fill=cyan!15, right=1.15cm of bdag] (sdag) {\textbf{Silver Processing DAG}\\(\texttt{silver\_processing\_dag})\\FileSensor + Validation\\1-to-1 Matching \& Energy};
\node[box, fill=green!12, below=0.35cm of sdag] (spg) {\textbf{Canonical PostgreSQL 16}\\\texttt{event, earthquake, source}\\\texttt{event\_group (1,274 rows)}\\\texttt{rejected\_record, state}};

% 4. Silver Parquet & PySpark Gold
\node[box, fill=green!12, right=1.15cm of sdag] (sparq) {\textbf{Silver Parquet Store}\\\texttt{datalake\_silver/}\\Uniform 13-Col Schema\\\texttt{TIMESTAMP(MICROS)}};
\node[box, fill=yellow!20, below=0.35cm of sparq] (gdag) {\textbf{Gold PySpark DAG}\\(\texttt{gold\_processing\_dag})\\\texttt{SparkSession local[*]}\\\texttt{spark.read.parquet()}};

% 5. Gold Outputs
\node[box, fill=orange!15!yellow!25, right=1.15cm of sparq] (ggov) {\textbf{Governance Gold Parquet}\\\texttt{datalake\_gold/}\\\texttt{governance\_*.parquet}\\(113 KPI Rows)};
\node[box, fill=purple!12, below=0.35cm of ggov] (gstory) {\textbf{Storytelling Gold Parquet}\\\texttt{datalake\_gold/}\\\texttt{storytelling\_*.parquet}\\(84 SAAI \& Trend Rows)};

\draw[arrow] (s1) -- (bdag);
\draw[arrow] (s2.east) -| ([xshift=-0.45cm]bdag.west) -- (bdag.west);
\draw[arrow] (bdag) -- (bstore);
\draw[arrow] (bstore.east) -| ([xshift=-0.45cm]sdag.west) -- (sdag.west);
\draw[arrow] (sdag) -- (spg);
\draw[arrow] (sdag) -- (sparq);
\draw[arrow] (sparq) -- (gdag);
\draw[arrow] (gdag.east) -| ([xshift=-0.45cm]ggov.west) -- (ggov.west);
\draw[arrow] (gdag) -- (gstory);

\begin{scope}[on background layer]
\node[layer, fit=(s1)(s2), label={[font=\scriptsize\bfseries\sffamily, text=navyblue]above:1. Seismic Networks}] {};
\node[layer, fit=(bdag)(bstore), label={[font=\scriptsize\bfseries\sffamily, text=navyblue]above:2. Bronze Layer}] {};
\node[layer, fit=(sdag)(spg), label={[font=\scriptsize\bfseries\sffamily, text=navyblue]above:3. Silver \& Canonical DB}] {};
\node[layer, fit=(sparq)(gdag), label={[font=\scriptsize\bfseries\sffamily, text=navyblue]above:4. PySpark Gold Engine}] {};
\node[layer, fit=(ggov)(gstory), label={[font=\scriptsize\bfseries\sffamily, text=navyblue]above:5. Gold Layer (Parquet)}] {};
\end{scope}
\end{tikzpicture}%
}
\caption{Updated three-layer Medallion and Canonical PostgreSQL architecture implemented in Workshop~\#3, integrating multi-source Bronze extraction, Silver validation/quarantine and cross-source matching, and PySpark \texttt{local[*]} Gold aggregations.}
\label{fig:w3_architecture}
\end{figure*}
